\section{Scope and separation of the mathematical layers}
The central sufficient-class theorem no longer requires an exact common finite primary row. Its assumptions are an actual acquired-mass envelope, a conditional task comparison characterized by Theorem~\ref{thm:fiber}, an all-strategy effort drift, and a verified finite-state acquisition program. These assumptions are not necessary for every finite-resource experiment. The fiber criterion is necessary and sufficient only for its specified all-prediction calibration mechanism. The reference acquired law is not a formal reachable support; alternative unknown-model laws are compared to it through actual conditional measure inequalities.

The theorem is load-bearing across the two raw families: their state-dependent primary laws and gate probabilities violate fixed-row closure, while the same effort, mass and calibration argument proves their risk curves. The declared finite physical programs, not their potentially infinite posterior trajectories, supply the implementation. The general Borel occupation theorem clarifies exact attainability of a terminal submeasure but does not reduce arbitrary hidden experiments to finite-state controllers.

Physical-command compatibility is invariant under stateless resource-preserving recodings. It does not compute the minimal automaton for an arbitrary protocol, and it is not preserved by adding a free phase-aware primitive or a free external clock. The two-port application has a matching live-state certificate. The numeral $M$ always counts a total alphabet for autonomous observers and a dictionary for checkpoint observers. Temporary workspace, program length, precision, call count, physical time and simulator state remain separate resources.

The regular marked geometric calculation is inherited mathematics, not a newly claimed quantization exponent. New conclusions concern its all-history predictable transfer, the criterion for its task calibration, unknown/history-dependent heralding, physical presentation control, and the actual atomic transition. The finite-bit example provides an explicit correct certificate and an interval, not an optimal arithmetic/compiler theorem. The same-task product theorem uses actual alternatives; general permitted errors produce only certified upper perturbations. No total-variation estimate is promoted to an LDP equivalence or to control of an uncertified mean stopping time.

A near-necessary acquisition classification for arbitrary unknown kernels, matching curves when the conditional mass envelopes fail, a full joint vanishing-$q$ law, efficient global program synthesis, and sharp lower bounds for program/workspace/physical time are not proved. Nor are unrestricted infinite-horizon or unbounded-operator assertions. A standard natural sensing problem with an independently established external open-problem significance is not claimed solved here. The finite-state program and the raw physical families are specified mathematical experiments, and novelty and significance require independent literature and proof assessment.

Historical companion articles are optional provenance and preserve distinct results on conditional pooling, continuation geometry, Bayesian curvature, finite-time allocation and stopped occupation. They are not unstated premises of this article. The main text contains the definitions, central proof, both raw verifications and the error composition needed for its own claims.
